% !TeX root = ../kazymyr-thesis.tex

\chapter{Related Work}\label{chapter:related-work}

This chapter is ordered by how close the work sits to this thesis, nearest
first. Section~\ref{sec:rw-simulators} covers simulation as a method for
studying concurrency control, and the base simulator this thesis extends,
which supplies the platform, the workload model and the two results the
evaluation starts from.
Section~\ref{sec:rw-mode-selection} covers the systems that already tag
each transaction with an execution mode at runtime and hand the decision to
an oracle, which is the mechanism used here as well. The work after those
is near on one axis each. Section~\ref{sec:rw-acid-base} covers running
ACID and BASE together, changing a consistency guarantee while the system
runs, and stating inconsistency as a number to be bounded.
Section~\ref{sec:rw-multimodel} covers the modelling framework the setting
comes from, and Section~\ref{sec:rw-gap} states what none of it reaches.

\section{Simulation of transactional systems}\label{sec:rw-simulators}

This section covers the simulators that this thesis builds its method on.
The first two are single-model studies that establish the method itself.
The last two are the multi-model simulator this thesis extends.

Transaction simulators have a long history in single-model settings. Ries
and Stonebraker~\cite{ries1979locking} used one to study locking
granularity under a workload held in equilibrium, and Dandamudi and
Au~\cite{dandamudi1991locking} extended that analysis to multiprocessor
database systems.

Both establish the method this thesis relies on: comparing
concurrency-control policies under identical synthetic workloads, in
virtual time, so that the comparison is not confounded by the host
machine. Neither concerns itself with the object of study here. Both
assume a single model and a single execution semantics, and the parameter
they vary, lock granularity or processor count, is a property of the
concurrency-control implementation, not a choice about how much
inconsistency the system is allowed to carry.

\paperbreak
Zacouris and Acosta~\cite{zacouris2025m2ts} carry that method into
multi-model systems. Their simulator models a transactional server in which
an update to one model triggers consistency-preserving updates in the
models that overlap it, executed by bookkeeper transactions. It runs under
ACID only, with strong strict two-phase locking, preclaiming and cautious
waiting, and it lets the user vary concurrency, transaction size, model
count, lock granularity and CPR structure. Its result is that throughput
collapses as the number of models grows, and that lock acquisition rather
than processing is the bottleneck, which is what motivates relaxing ACID at
all.

The subsequent study~\cite{zacouris2026mxs} does that. It adds BASE, in
which a transaction takes no locks, conflicting writes are settled by last
write wins, and the consistency-preserving updates are deferred to a
reconciliation queue served by dedicated reconciler processors. It adds
semantic overlap as a parameter and the inconsistency metric this thesis
controls, that is the percentage of resources awaiting repair, reported at
the median, the 95th and the 99th percentile. Its comparison is
where this thesis starts. Under ACID the consistency-preserving work runs
inside the triggering transaction, and throughput degrades sharply as
semantic overlap and CPR density grow. Under BASE it is deferred, and
foreground throughput becomes largely insensitive to both, at the price of
transient inconsistency whose level depends on the processors allocated to
reconciliation.

That work provides the platform, the workload model and two results that
are the direct starting point here. The first is the
throughput/consistency trade-off governed by reconciler provisioning: too
few reconcilers leave inconsistency high, too many starve foreground
execution. Section~\ref{sec:x1-validation} reproduces its shape and uses
it to validate the extended simulator. The second is that the best
allocation depends on the configuration of the multi-model system, which
is what makes reconciler provisioning an awkward control: a knob whose
correct setting must be found again for every deployment.

Its stated assumptions are inherited here as well, and they are the reason
several questions are out of scope in this thesis. Execution is fully
centralised, so availability is not modelled, there is no fault tolerance
and no network delay. The system is of constant size, since writes change
resources but never create or destroy them. The workload holds a constant
number of transactions in flight, and a terminal-based workload is named
there as future work. That terminal workload is the \texttt{TERMINAL}
profile of Section~\ref{sec:prelim-workload-profiles}, which was still
being built while these experiments ran and is the reason the adaptive
oracle is tested only under a stationary load
(Section~\ref{sec:approach-stationary-load}).

What that work does not provide is a way to select the execution semantics
\emph{per transaction} and \emph{at runtime}. There, ACID and BASE are
alternative configurations of a whole run. Here they coexist within one
run, which requires the arbiter of Section~\ref{sec:approach-arbiter} and
makes the tolerated inconsistency a configurable parameter rather than an
outcome of the provisioning.

\section{Choosing an execution mode at runtime}\label{sec:rw-mode-selection}

Closest after the base simulator is a line of work that already does what
this thesis does mechanically. It tags each transaction with one of two
execution modes while the system runs, and routes it accordingly. It is
kept apart from the sections that follow because the choice it makes is
never a choice about consistency.

Hybrid Transactional Replication~\cite{kobus2018htr} copies a service onto
several machines and runs each transaction in one of two modes. In
\emph{deferred-update} mode a single replica runs the transaction on its own
copy without locking anything and tells the other replicas what it read and
wrote only at the end, so they may reject it and force a retry. In
\emph{state-machine} mode every replica runs the same transaction on its own
copy in one agreed order, so nothing can conflict and no transaction is ever
aborted. Deferred update therefore buys parallelism at the risk of wasted
work, and the state machine buys freedom from aborts at the cost of doing
the same work on every replica and in sequence. Which mode a transaction
gets is decided at runtime by a component the authors also call an
\emph{oracle}, and their oracle learns from the behaviour of past
transactions. PolyCert~\cite{couceiro2011polycert} makes the same kind of
per-transaction choice between certification protocols in a replicated
in-memory store, using a learned model. MorphR~\cite{couceiro2013morphr}
switches the replication protocol of the whole system rather than of one
transaction, and contributes the machinery for changing protocol safely
while work is in flight.

Together these establish that a per-transaction mode tag issued at runtime
is a workable design, and that a narrow oracle interface is the right way
to keep the policy separate from the execution paths, which is the shape
adopted in Section~\ref{sec:approach-oracle-interface}. What they do not
contribute is the thing being traded. All of their modes deliver the same
consistency, so their oracles optimise throughput and latency alone and no
operator can ask for a consistency level at all. There is no dirty state to
read, no inconsistency to bound, and therefore no equivalent of the mixed
state that Section~\ref{sec:approach-arbiter} has to make well defined.

\section{Combining ACID and BASE}\label{sec:rw-acid-base}

The three subsections below take one axis each: running the two semantics
side by side, changing a guarantee while the system runs, and stating
inconsistency as a number that can be bounded.

\subsection{Static classification of transactions}\label{sec:rw-static}

Salt~\cite{xie2014salt} is the nearest work on running ACID and BASE side
by side in one system. A developer picks the few performance-critical ACID
transactions of an application and rewrites them as BASE transactions, each
split into smaller steps. Every other transaction stays ACID. What keeps
that mixture safe is a rule the authors call Salt Isolation. BASE
transactions are allowed to see each other's unfinished work at defined
points, which is where their speed comes from, while ACID transactions are
held apart from that unfinished work and get the same isolation they would
have had in a system with no BASE transactions in it.

What Salt provides for this thesis is the argument that the two semantics
can coexist in one system without collapsing into the weaker one, together
with the precedent for the rule that achieves it. Keeping ACID transactions
away from BASE work that is still in flight is also what the claim table of
Section~\ref{sec:approach-claim-table} does here.

What Salt does not provide is the decision this thesis is about. It fixes
\emph{which} transactions are BASE statically, in the application source,
whereas the question here is \emph{when} to use each mode, at runtime, from
measured system state. Salt also has no background repair work, so no
resource is ever left carrying an update that still has to be applied, and
the direction in which ACID waits for reconciliation has no counterpart
there. Nor does the setting match. The prototype is one relational schema
split over the partitions of a distributed database, so the inconsistency
that automatic propagation between overlapping models creates, which is the
only kind in this thesis, cannot arise at all.

\paperbreak
RedBlue consistency~\cite{li2012redblue} makes a similar classification at
a finer granularity: operations labelled blue execute without coordination
and operations labelled red require it, with the labelling fixed at design
time and constrained by commutativity requirements.

It contributes the finer unit of classification, an operation rather
than a transaction, and a criterion for when weak execution is safe at
all. It does not contribute a runtime element: the labels are a property
of the program, not of the current state of the system, and the
commutativity requirement that justifies them has no counterpart here,
where the deferred work is a consistency-preserving update whose effect is
visible in the dirty fraction rather than a commuting operation.

% Closes the subsection: what the two entries above have in common.
\paperbreak
Both classifications are applied by hand. In Salt a developer decides which
transactions are worth rewriting as BASE ones and rewrites them. In RedBlue
a programmer splits every operation into its execution paths and then
labels each path, which first requires the invariants of the application to
be written down, because an operation has to be red if it fails to commute
with another or if it could break one of those invariants. The marking is a
property of the source code in both cases. It is made once, before the
system runs, by someone reasoning about the whole application, and it cannot
answer to anything the system does afterwards. Removing that last
limitation is what the mode oracle of Chapter~\ref{chapter:approach} is
for.

\subsection{Runtime-adaptive consistency}\label{sec:rw-adaptive}

Consistency Rationing~\cite{kraska2009rationing} is the closest
methodological precedent. Data is divided into categories that receive
different guarantees, and the middle category switches its consistency
level at runtime according to a policy driven by monitored statistics and
expected cost.

What it gives this thesis is the shape of the mechanism: monitor
something, compare it against a policy, and change the guarantee while the
system is running. The mode oracle of Chapter~\ref{chapter:approach} is
the same idea applied to a different signal. What it does not give is the
signal itself. Rationing steers by an expected monetary cost, so the
consistency that results is whatever the cost model happens to produce,
and the categories are assigned to data items rather than decided per
transaction at admission.

\paperbreak
OptCon~\cite{sidhanta2016optcon} lets the operator name the target instead
of a proxy for it. A user gives a staleness bound and a latency bound in a
service level agreement, and a learned model predicts, for each operation,
which consistency setting of the underlying quorum store will meet them.

This is the nearest precedent for the framing of RQ2, that an operator
states a level and the system delivers it, and OptCon takes its setting in
the units of the quantity it bounds, as $c$ does. What it does not give is
the unit of the decision. OptCon chooses how a single read is served, not
whether a whole transaction and the repairs it triggers run atomically. It
is also best-effort against the stated bound rather than measured against
it, and it has no notion of work owed by one model to another.

\subsection{Bounded inconsistency}\label{sec:rw-bounded}

Epsilon-serializability~\cite{pu1991esr} is the oldest statement of the
idea this thesis depends on. A transaction that can tolerate some
inconsistency is allowed to run non-serialisably, up to a bound
$\varepsilon$ that is stated in advance, and it gains concurrency in
exchange. The bound is held by \emph{divergence control}, a family of
algorithms that stands to the bound as concurrency control stands to
serialisability. Wu, Yu and Pu~\cite{wu1992divergence} give versions of it
built on two-phase locking, on timestamp ordering and on optimistic
execution.

This is the closest ancestor of the critical threshold, and closer than
the replication work below, because it puts a numerical inconsistency
budget inside a transactional system and enforces it with the locking
machinery that system already has. What it does not offer is the shape of
the budget. Its bound is per transaction and concerns how far the values
one transaction reads may be from the current ones, whereas $c$ is a
single bound over the whole resource population, and the work it bounds is
repair owed by one model to another rather than reads taken out of order.
Nothing in that line adapts the bound, or the mode, while the system runs.

\paperbreak
Yu and Vahdat~\cite{yu2000conit} treat consistency as a continuous
quantity and allow applications to bound the divergence they tolerate
along several axes through conits.

This supplies the vocabulary used throughout this thesis: inconsistency as
a number that an application may bound, rather than a property that is
either held or lost. The critical dirty fraction is a conit-style
numerical bound, and the tracking result of
Section~\ref{sec:x2-critical-sweep} is evidence that it behaves as one
here. What conits do not supply is a way to hold such a bound in a
multi-model system. The enforcement mechanisms are those of replicated
state, that is pushing writes between replicas to keep divergence within
the declared limits, and they have no equivalent for inconsistency created
by consistency-preserving updates between models.

\section{Multi-model consistency preservation}\label{sec:rw-multimodel}

Vitruvius~\cite{klare2021vitruvius} and the line of work on consistency
preservation rules provide the setting itself: automated propagation of
updates between models that represent the same system from different
viewpoints.

The name stands for view-centric engineering using a virtual underlying
single model. Rather than merging the models of a project into one
redundancy-free model, Vitruvius keeps them as they are and makes them
behave as if they were one, by writing the rules that tie them together
down explicitly. Engineers work through views, and a change made in a view
is propagated by those rules into every other model that shares the
affected information.

It defines what a consistency-preserving rule is, how such rules are
specified, and how they are executed when a model changes. That is the
machinery which generates the workload studied here. It is also where the
numbers come from: the base simulator is calibrated with measurements taken
from multi-model systems built in Vitruvius~\cite{zacouris2026mxs}, so the
parameter values of Chapter~\ref{chapter:experiments} trace back to it.
What it does not do is treat the propagation as something with a cost or a
delay. There are no transactions, no concurrent access to arbitrate and no
option to defer a rule. Propagation is a correctness mechanism to be
specified and executed, not a source of transient inconsistency to be
measured and bounded, so questions of scheduling, deferral and the
resulting dirty fraction fall outside it.

\section{Discussion and delimitation}\label{sec:rw-gap}

Three lines of delimitation run through the work above. It is worth being
precise about the first, because the obvious way to state it would claim
too much. Tagging each transaction with an execution mode at runtime is not
new. Hybrid Transactional Replication and PolyCert already do it, and they
do it with a learned policy rather than a threshold. Against those systems
the difference is not the mechanism but what the two modes are. Their modes
differ in how a transaction is executed and agree on what it guarantees, so
there is one quantity to optimise and none to bound. Here the two modes
differ in what they guarantee, which is what creates both the inconsistency
to be bounded and the mixed state that
Section~\ref{sec:approach-arbiter} has to define. Against the base
simulator, where ACID and BASE do differ in what they guarantee, the
difference is granularity and timing instead: the semantics are chosen per
transaction at admission rather than per run at configuration time.

Against the systems that do adapt consistency at runtime, the difference is
what the operator actually sets. Consistency Rationing is steered by an
expected cost and the base simulator by a reconciler count. Each is a
proxy: the operator sets one quantity and observes another. OptCon does
take a target in the right units, but it aims that target at how one read
is served, not at whether a transaction and the repairs it triggers run
atomically. Here $c$ is stated in the units of
the quantity it bounds and applies to the whole run, and the RQ3 result,
not the literature, is what supports that.

Beyond the base simulator, all of this work targets replicated or
partitioned storage under a single schema, where inconsistency means
divergence between copies of a datum. None of it addresses the
inconsistency that consistency-preserving updates create between
overlapping models. That is the gap this thesis fills.
